\section{Conclusions}
\label{sec:conclusions}

We have shown, constructively and reproducibly, that a single deterministic elastic
substrate---a 4D cubic brane lattice whose only nonlinearity is the Pythagorean link
length---can carry emergent low-energy \emph{symmetries} of two distinct kinds. First, the
Lorentzian signature of the effective metric is the sign pattern of the directional prestress,
and the timelike link's continuum limit is a genuine kinetic term with emergent inertia; the
result is a massless, per-sector Lorentz-invariant light cone with a closed-form calibration.
Second, the fluctuation operator about a finite-amplitude periodic carrier possesses an isolated
rank-3 complex eigenbundle whose gauge-invariant Wilczek--Zee curvature generates the full Lie
algebra $\su{3}$---generically, at leading order in the carrier amplitude, and satisfying the
non-abelian Bianchi identity---with the abelian trace as a candidate electromagnetic $U(1)$.
Emergent local gauge invariance then forces the leading dynamics of both sectors to Maxwell and
Yang--Mills form, with substrate-induced positive couplings and a quantitatively controlled
$U(3)$ (not $U(4)$) rank selection.

The central methodological commitment is honesty about scope: the color-carrying background is a
stable but higher-energy texture rather than the ground state; the emergent Lorentz invariance is
approximate, and that very approximateness is what lets the emergent gauge bosons evade the
Weinberg--Witten theorem; and the matter sector is deferred. Within those boundaries, the work is
a concrete demonstration that emergent $U(1)\times SU(3)$ gauge structure and Lorentzian geometry
can descend from the same carrier geometry of one deterministic lattice.

The natural next steps follow from the stated limitations: constructing self-confined solitons
(mass, color charge, and the coupling to the gauge field derived here); settling the selection and
energetics of the vacuum carrier; and testing the dual-observer route to full Lorentz universality
at the matter scale. Each is a well-posed problem built directly on the interfaces established
here.
